\par\SourceLeaf{528}Die Vorschrift lässt sich also auch so ausdrücken; (stets unter der Voraussetzung $p\equiv5\ (\mathrm{mod}.8)$)
\[
\begin{aligned}
\text{Es ist }b&\equiv a+1\ (\mathrm{mod}.8),&\quad&\text{wenn }m\text{ gerade}\\
b&\equiv a+5&&\text{wenn }m\text{ ungerade.}
\end{aligned}
\]
Ich wage noch keine Vermutung darüber, ob ein noch einfacheres Kriterium möglich ist, mit dem sich der Fall des geraden $m$ von dem ungeraden Fall \textit{im Voraus} unterscheiden ließe, also ohne den Wert von $m$ selbst zu kennen, da diese Verbindung, wie ich schon oben bemerkt habe, noch ganz neu ist.
\newpage
\par Für den Fall $p\equiv1\ (\mathrm{mod}.8)$ bleibt zwar die obige Kongruenz $b\equiv2k+a-1\ (\mathrm{mod}.8)$ richtig, sie entscheidet aber nicht mehr über das Vorzeichen von $b$, da sowohl der positive als auch der negative Wert von $b$ ihr genügt. Hier ist nämlich $k$ immer gerade, $=2m$ (wenn $m$ dieselbe Bedeutung wie oben hat) oder $=2m+2$, wenn $m$ die Anzahl der Zerlegungen von $p$ in 3 positive \textit{ungleiche} Quadrate bezeichnet, und $b\equiv0\ (\mathrm{mod}.4)$ oder $b\equiv-b\ (\mathrm{mod}.8)$. Ich vermute, dass der Fall $p\equiv1\ (\mathrm{mod}.8)$ oder $b\equiv0\ (\mathrm{mod}.4)$ einer tieferen Untersuchung bedarf und vielleicht wiederum
\[
\begin{aligned}
b&\equiv4\ (\mathrm{mod}.8)&&\text{leichter als }b\equiv0\ (\mathrm{mod}.8)\\
b&\equiv8\ (\mathrm{mod}.16)&&\text{leichter als }b\equiv0\ (\mathrm{mod}.16)
\end{aligned}
\]
usw.
\par\Needspace{5\baselineskip}Mit ausgezeichneter Hochachtung beharre ich
\par\nobreak\noindent\hfill Ihr freundschaftlich ergebenster
\par\nobreak\noindent Göttingen, den 30. Mai 1828.\hfill C. F. \textsc{Gauss}.
\BookRule{.25\textwidth}
